% GENERATED FILE. Edit canonical lessons, then run npm run generate:books.
% canonical-source-hash: fnv1a64:33a11ae96bc17ea6
% canonical-lessons: TA-C196-kili, TA-C196-ila, TA-C196-aluttu, TA-C196-unar, TA-C196-ura-vai

\chapter{To tear, To lose, To press, To feel, To soak}
\label{ch:ta-to-tear-to-lose-to-press-to-feel-to-soak}
\hlchaptermodality{196}

\begin{chapteropening}
\textbf{By the end of this chapter:} \emph{I can say to tear, to lose, to press, to feel and to soak.}

Complete the last lesson of chapter 196: I can say to tear, to lose, to press, to feel and to soak.
\end{chapteropening}

\section[\texorpdfstring{\ta{கிழி} (kiḻi)}{கிழி (kiḻi)}]{\ta{கிழி} (kiḻi) — to tear}

\label{lesson:TA-C196-kili}



\begin{quote}
Before the new one: say the Tamil for to use, then the Tamil for to try.
\end{quote}

\subsection*{You'll want to know: \ta{கிழி}}
\textbf{\ta{கிழி}} — \emph{kiḻi} — \textquotedblleft{}to tear\textquotedblright{}.

\begin{grammarlens}[title={What you've built}]
One more word for this chapter.
\end{grammarlens}

\subsection*{Your turn}
\begin{itemize}
\raggedright
  \item \emph{Say it:} \emph{kiḻi}
  \item \emph{Say it:} \emph{kiḻi}, once more
  \item \emph{Say it:} say \emph{kiḻi}
  \item \emph{Recall:} say the Tamil for to use, then the Tamil for to try, then say \emph{kiḻi} again
\end{itemize}

\begin{tcolorbox}[breakable,colback=teal!4,colframe=teal!35!black,title={Before you move on}]
What does \ta{கிழி} mean? (\textquotedblleft{}To tear\textquotedblright{}.) Say it once more.
\end{tcolorbox}
\section[\texorpdfstring{\ta{இழ} (iḻa)}{இழ (iḻa)}]{\ta{இழ} (iḻa) — to lose}

\label{lesson:TA-C196-ila}



\begin{quote}
Before the new one: say the Tamil for to try, then the Tamil for to tear.
\end{quote}

\subsection*{You'll want to know: \ta{இழ}}
\textbf{\ta{இழ}} — \emph{iḻa} — \textquotedblleft{}to lose\textquotedblright{}.

\begin{grammarlens}[title={What you've built}]
One more word for this chapter.
\end{grammarlens}

\subsection*{Your turn}
\begin{itemize}
\raggedright
  \item \emph{Say it:} \emph{iḻa}
  \item \emph{Say it:} \emph{iḻa}, once more
  \item \emph{Say it:} say \emph{iḻa}
  \item \emph{Recall:} say the Tamil for to try, then the Tamil for to tear, then say \emph{iḻa} again
\end{itemize}

\begin{tcolorbox}[breakable,colback=teal!4,colframe=teal!35!black,title={Before you move on}]
What does \ta{இழ} mean? (\textquotedblleft{}To lose\textquotedblright{}.) Say it once more.
\end{tcolorbox}
\section[\texorpdfstring{\ta{அழுத்து} (aḻuttu)}{அழுத்து (aḻuttu)}]{\ta{அழுத்து} (aḻuttu) — to press}

\label{lesson:TA-C196-aluttu}



\begin{quote}
Before the new one: say the Tamil for to tear, then the Tamil for to lose.
\end{quote}

\subsection*{You'll want to know: \ta{அழுத்து}}
\textbf{\ta{அழுத்து}} — \emph{aḻuttu} — \textquotedblleft{}to press\textquotedblright{}.

\begin{grammarlens}[title={What you've built}]
One more word for this chapter.
\end{grammarlens}

\subsection*{Your turn}
\begin{itemize}
\raggedright
  \item \emph{Say it:} \emph{aḻuttu}
  \item \emph{Say it:} \emph{aḻuttu}, once more
  \item \emph{Say it:} say \emph{aḻuttu}
  \item \emph{Recall:} say the Tamil for to tear, then the Tamil for to lose, then say \emph{aḻuttu} again
\end{itemize}

\begin{tcolorbox}[breakable,colback=teal!4,colframe=teal!35!black,title={Before you move on}]
What does \ta{அழுத்து} mean? (\textquotedblleft{}To press\textquotedblright{}.) Say it once more.
\end{tcolorbox}
\section[\texorpdfstring{\ta{உணர்} (uṇar)}{உணர் (uṇar)}]{\ta{உணர்} (uṇar) — to feel}

\label{lesson:TA-C196-unar}



\begin{quote}
Before the new one: say the Tamil for to lose, then the Tamil for to press.
\end{quote}

\subsection*{You'll want to know: \ta{உணர்}}
\textbf{\ta{உணர்}} — \emph{uṇar} — \textquotedblleft{}to feel\textquotedblright{}.

\begin{grammarlens}[title={What you've built}]
One more word for this chapter.
\end{grammarlens}

\subsection*{Your turn}
\begin{itemize}
\raggedright
  \item \emph{Say it:} \emph{uṇar}
  \item \emph{Say it:} \emph{uṇar}, once more
  \item \emph{Say it:} say \emph{uṇar}
  \item \emph{Recall:} say the Tamil for to lose, then the Tamil for to press, then say \emph{uṇar} again
\end{itemize}

\begin{tcolorbox}[breakable,colback=teal!4,colframe=teal!35!black,title={Before you move on}]
What does \ta{உணர்} mean? (\textquotedblleft{}To feel\textquotedblright{}.) Say it once more.
\end{tcolorbox}
\section[\texorpdfstring{\ta{ஊற} \ta{வை} (ūṟa vai)}{ஊற வை (ūṟa vai)}]{\ta{ஊற} \ta{வை} (ūṟa vai) — to soak}

\label{lesson:TA-C196-ura-vai}



\begin{quote}
Before the new one: say the Tamil for to press, then the Tamil for to feel.
\end{quote}

\subsection*{You'll want to know: \ta{ஊற} \ta{வை}}
\textbf{\ta{ஊற} \ta{வை}} — \emph{ūṟa vai} — \textquotedblleft{}to soak\textquotedblright{}.

\begin{grammarlens}[title={What you've built}]
One more word for this chapter.
\end{grammarlens}

\subsection*{Your turn}
\begin{itemize}
\raggedright
  \item \emph{Say it:} \emph{ūṟa vai}
  \item \emph{Say it:} \emph{ūṟa vai}, once more
  \item \emph{Say it:} say \emph{ūṟa vai}
  \item \emph{Recall:} say the Tamil for to press, then the Tamil for to feel, then say \emph{ūṟa vai} again
\end{itemize}

\begin{tcolorbox}[breakable,colback=teal!4,colframe=teal!35!black,title={Before you move on}]
What does \ta{ஊற} \ta{வை} mean? (\textquotedblleft{}To soak\textquotedblright{}.) Say it once more.
\end{tcolorbox}
